\documentclass{article}
\usepackage[utf8]{inputenc}
\usepackage[margin=2.5cm]{geometry}
\usepackage{hyperref}

% Formular-Einstellungen global festlegen
\hypersetup{
	pdftitle={Mein großes Formular},
	pdfpagemode=UseNone,
	pdfstartview=FitH
}

\newcommand{\dreierwahl}[2]{% #1=Name, #2=Fragetext
	#2 & \ChoiceMenu[radio, name=#1, radiosymbol=\ding{108}]{}{ja=J} 
	& \ChoiceMenu[radio, name=#1, radiosymbol=\ding{108}]{}{nein=N} 
	& \ChoiceMenu[radio, name=#1, radiosymbol=\ding{108}]{}{vielleicht=V} \\
}

\begin{document}
	
	\section*{Anamnesebogen (Beispiel f\"ur 100+ Felder)}
	
	\begin{Form} % Wichtig: Alles muss innerhalb der Form-Umgebung stehen
\noindent % Verhindert Einrückung der Tabelle
\begin{tabular}{@{}p{4cm} l}
	\textbf{Name:}         & \TextField[name=name, width=8cm, height=1cm, charsize=14pt]{} \\[2ex]
	\textbf{Vorname (hoch):} & \TextField[name=vorname, width=8cm, height=1cm, charsize=14pt]{} \\[5ex]
	\textbf{Geburt:}       & \TextField[name=geburt, width=4cm, height=0.5cm, bordercolor={0.8 0.8 0.8}]{} \\
\end{tabular}

\begin{description}
	\item[Name:] \TextField[name=name, width=8cm]{}
	\item[Vorname:] \TextField[name=vorname, width=8cm]{}
\end{description}
		
		
		\subsection*{Gesundheitsfragen}
		
		% Hier kannst du mit einer einfachen LaTeX-Schleife oder Copy-Paste arbeiten
		\noindent
		Sind Sie allergisch gegen: \\
		\CheckBox[name=all_pollen]{Pollen} \hspace{1cm}
		\CheckBox[name=all_staub]{Hausstaub} \hspace{1cm}
		\CheckBox[name=all_tier]{Tierhaare}
		
		\noindent
		\textbf{Wie zufrieden bist du mit der L\"osung?} \\
		\ChoiceMenu[radio, name=Zufriedenheit, radiosymbol=\ding{108}]{}{Sehr gut=1, Gut=2, Mittel=3, Schlecht=4}
		
		\vspace{1cm}
		
		\begin{tabular}{lccc}
			\textbf{Frage} & \textbf{Ja} & \textbf{Nein} & \textbf{Vielleicht} \\ \hline
			\dreierwahl{q1}{Gef\"allt dir dieser Ansatz?}
			\dreierwahl{q2}{Soll es so weitergehen?}
			\dreierwahl{q3}{Sind 100 Felder viel?}
		\end{tabular}
		
		\vspace{0.5cm}
		
		\noindent
		\textbf{Zusatzbemerkungen (mehrzeilig):} \\
		\TextField[name=bemerkung, multiline=true, width=\textwidth, height=3cm, bordercolor={0.8 0.8 0.8}]{}
		
	\end{Form}
	
\end{document}